\section{第 12 章：规范化}

\begin{taplauthorsolutionentrybox}
\phantomsection
\label{sol:author:12.1.1}
\textbf{12.1.1 解答}

问题出在应用情形。设我们要证明 \(t_1\,t_2\) 可规范化。由归纳假设，
\(t_1\) 与 \(t_2\) 都可规范化；记其范式分别为 \(v_1\) 与 \(v_2\)。
由\taplnamedref{lem:9.3.1}{类型关系的反演引理}，\(v_1\) 具有某个函数类型
\(T_{11}\to T_{12}\)。再由典范形式引理（\cref{lem:9.3.4}），\(v_1\)
必为 \(\lambda x:T_{11}.\,t_{12}\)。

可是，把 \(t_1\,t_2\) 归约到
\((\lambda x:T_{11}.\,t_{12})\,v_2\) 并没有得到范式，因为此时
\rulename{E-AppAbs} 仍然适用，产生
\([x\mapsto v_2]t_{12}\)。要完成证明，就必须说明这个新项也可规范化；
而结构归纳在这里无法继续，因为替换可能按照 \(x\) 在 \(t_{12}\) 中出现的
次数复制 \(v_2\)，使新项比原项 \(t_1\,t_2\) 更大。

\taplsolutionbacklink{ex:12.1.1}{返回练习 12.1.1}
\end{taplauthorsolutionentrybox}

\begin{taplauthorsolutionentrybox}
\phantomsection
\label{sol:author:12.1.7}
\textbf{12.1.7 解答}

在定义 12.1.2 中增加两项：
\[
\begin{aligned}
\mathcal R_{\tapltype{Bool}}(t)
  &\quad\Longleftrightarrow\quad
    t\text{ 的求值会在有限步内终止},\\
\mathcal R_{T_1\times T_2}(t)
  &\quad\Longleftrightarrow\quad
    \begin{aligned}[t]
    &t\text{ 的求值会在有限步内终止，且}\\
    &\mathcal R_{T_1}(t.1)\text{ 与 }\mathcal R_{T_2}(t.2)
    \end{aligned}
\end{aligned}
\]

引理 12.1.4 的证明需要增加一个情形。设 \(T=T_1\times T_2\)。对于
“仅当”方向，假设 \(\mathcal R_T(t)\) 且 \(t\trans t'\)，需证
\(\mathcal R_T(t')\)。由 \(\mathcal R_T\) 的定义可知
\(\mathcal R_{T_1}(t.1)\) 与 \(\mathcal R_{T_2}(t.2)\)。规则
\rulename{E-Proj1} 和 \rulename{E-Proj2} 给出
\(t.1\trans t'.1\) 与 \(t.2\trans t'.2\)；对二者使用归纳假设，得到
\(\mathcal R_{T_1}(t'.1)\) 与 \(\mathcal R_{T_2}(t'.2)\)，再由定义
得到 \(\mathcal R_T(t')\)。
“如果”方向相同。

最后，要为引理 12.1.5 的证明补上各条新类型规则的情形。

\medskip\noindent
\textbf{情形 \rulename{T-If}：}

此时
\[
  t=\taplif{t_1}{t_2}{t_3},
  \quad
  \Gamma\vdash t_1:\tapltype{Bool},
  \quad
  \Gamma\vdash t_2:T,
  \quad
  \Gamma\vdash t_3:T
\]
其中 \(\Gamma=x_1:T_1,\ldots,x_n:T_n\)。令
\(\sigma=[x_1\mapsto v_1]\cdots[x_n\mapsto v_n]\)。由归纳假设，
\(\mathcal R_{\tapltype{Bool}}(\sigma t_1)\)、
\(\mathcal R_T(\sigma t_2)\) 和 \(\mathcal R_T(\sigma t_3)\)。
第一项说明 \(\sigma t_1\) 的求值会在有限步内终止；它又具有布尔类型，
所以其最终值只能是 \(\tapltrue\) 或 \(\taplfalse\)。因此条件表达式
\(\sigma t\) 最终会选择 \(\sigma t_2\) 或 \(\sigma t_3\)。两者均已满足
\(\mathcal R_T\)，再由引理 12.1.4 对求值保持 \(\mathcal R_T\) 的结论，
得到 \(\mathcal R_T(\sigma t)\)。

\medskip\noindent
\textbf{情形 \rulename{T-True} 与 \rulename{T-False}：}

立即成立。

\medskip\noindent
\textbf{情形 \rulename{T-Pair}：}

此时 \(t=\{t_1,t_2\}\)，且 \(T=T_1\times T_2\)。由归纳假设，
对 \(i=1,2\) 都有 \(\mathcal R_{T_i}(\sigma t_i)\)，所以两个分量的
求值都会在有限步内终止；按值求值因而也会把
\(\{\sigma t_1,\sigma t_2\}\) 求值为一个值。记 \(\sigma t_i\) 的范式
为 \(w_i\)；由引理 12.1.4，\(\mathcal R_{T_i}(w_i)\)。投影的求值规则给出
\[
  \{\sigma t_1,\sigma t_2\}.i\transmulti w_i
\]
再次使用引理 12.1.4，可得
\(\mathcal R_{T_i}(\{\sigma t_1,\sigma t_2\}.i)\)。于是由定义得到
\(\mathcal R_{T_1\times T_2}(\{\sigma t_1,\sigma t_2\})\)，也就是
\(\mathcal R_{T_1\times T_2}(\sigma t)\)。

\medskip\noindent
\textbf{情形 \rulename{T-Proj1}：}

由 \(\mathcal R_{T_1\times T_2}\) 的定义立即得到。

\medskip\noindent
\textbf{情形 \rulename{T-Proj2}：}

类似。

\taplsolutionbacklink{ex:12.1.7}{返回练习 12.1.7}
\end{taplauthorsolutionentrybox}
